\documentclass[12pt]{article}
\usepackage[T1]{fontenc}
\usepackage[utf8]{inputenc}
\usepackage{lmodern}
\usepackage{textcomp}
\usepackage{enumitem}
\usepackage{geometry}
\geometry{margin=1in}
\begin{document}
\begin{center}
Answer Key - History - Graduate Level Assessment 10 \
\small (Answers are ordered exactly as in the question paper.)
\end{center}

\section*{Multiple Choice}
\begin{enumerate}[label=\arabic*.]
\item \textbf{Answer:} D
\\ \textit{Options:} A) geological layers and sediment types.; B) the architectural style and layout of structures on the site.; C) pumice and evidence of pyroclastic surges.; D) the faunal assemblage from a site.; E) an osteological comparative collection from a site.
\\ \textit{Note:} The correct answer is based on the fact that the faunal assemblage, which includes the bones and remains of animals found at an archaeological site, provides crucial insights into the diet, hunting practices, and resource utilization of past societies, thus informing us about their subsistence strategies.
\item \textbf{Answer:} A
\\ \textit{Options:} A) Caral; B) Chichen Itza; C) The Parthenon; D) Machu Picchu; E) Angkor Wat
\\ \textit{Note:} Caral, an ancient city in Peru, is recognized for its monumental earthworks and architectural features, including concentric ridges in a semicircular layout, which date back approximately 3,500 years, highlighting its significance in pre-Columbian history.
\item \textbf{Answer:} C
\\ \textit{Options:} A) understanding wave patterns; B) cloud patterns; C) magnetic orientation to the poles; D) bird flight paths; E) currents and wind patterns
\\ \textit{Note:} The native navigators of the Pacific relied on celestial navigation, ocean currents, wind patterns, and other natural indicators rather than magnetic orientation to the poles, which was not a technique they utilized.
\item \textbf{Answer:} A
\\ \textit{Options:} A) Sunda; Sahul; B) Sunda; Wallacea; C) Gondwana; Beringia; D) Sunda; Gondwana; E) Beringia; Sunda
\\ \textit{Note:} The correct answer is based on the geological understanding that during the Pleistocene era, the islands of Java, Sumatra, Bali, and Borneo were part of the continental shelf known as Sunda, while the landmass connecting Australia, New Guinea, and Tasmania is referred to as Sahul, reflecting the distinct biogeographical divisions of the region.
\item \textbf{Answer:} E
\\ \textit{Options:} A) They had no domesticated animals but hunted wild animals for meat.; B) They had domesticated animals, but they did not consume their meat.; C) They bred and raised a specific species of deer for food.; D) They were vegetarian and did not consume meat.; E) They brought animals to the site from outside of the region.
\\ \textit{Note:} The analysis of animal teeth revealed that the isotopic signatures were inconsistent with local fauna, indicating that the builders of Stonehenge imported animals from other regions for their construction and ceremonial purposes.
\end{enumerate}

\section*{True / False}
\begin{enumerate}[label=\arabic*.]
\setcounter{enumi}{5}
\item \textbf{Answer:} True
\\ \textit{Options:} A) True; B) False
\\ \textit{Note:} According to the passage: Obsessed.
\item \textbf{Answer:} False
\\ \textit{Options:} A) True; B) False
\\ \textit{Note:} The correct answer is: Serbia
\item \textbf{Answer:} False
\\ \textit{Options:} A) True; B) False
\\ \textit{Note:} The correct answer is: 1569
\item \textbf{Answer:} True
\\ \textit{Options:} A) True; B) False
\\ \textit{Note:} According to the passage: John Williams
\item \textbf{Answer:} False
\\ \textit{Options:} A) True; B) False
\\ \textit{Note:} The correct answer is: two towers
\end{enumerate}

\section*{Long Form Response}
\begin{enumerate}[label=\arabic*.]
\setcounter{enumi}{10}
\item \textbf{Answer:} Dr. T.R.M. Howard
\\ \textit{Note:} Expected answer: Dr. T.R.M. Howard
\item \textbf{Answer:} turned over control of his armies to the state
\\ \textit{Note:} Expected answer: turned over control of his armies to the state
\end{enumerate}

\vfill
\noindent\textit{End of Answer Key}
\end{document}